% MT039 / explanation — canonical editable content.
仍以\(z=yf(x)\)为横轴。负值分错，正值分对；损失曲线的高度表示当前付出多少代价，斜率表示把分数挪一点能改变多少代价。两者要分开看。


\textbf{右边\(z\ge1\)时损失为0。} 分对且达到足够间隔的样本不再推动数据损失下降，这与hinge的右侧相似。模型能把注意力放到还没有满足间隔的点上。


\textbf{最左边是一个较高但逐渐变平的平台。} 极端错分点仍有损失，绝不是越错损失越接近0。但局部斜率趋近0，继续改变它的分数只能带来很小的损失变化。因此梯度法从它收到的更新信号很弱。


这有两面作用：若极端点是异常值，弱化影响可避免模型为它牺牲很多正常点；若它是有用但目前分错的样本，过小的梯度也可能让模型难以把它救回来。不能只写“稳健”而忽略训练困难。


图形还是非凸的：凸函数的斜率从左往右不能变小，而本图从左端近乎水平，逐渐变为更陡的下降段，斜率由接近0变得更负。因此不能直接使用凸SVM“局部最优即全局最优”的结论。原解中“远端错分的处罚下降”应准确读成“梯度影响下降”，平台高度并没有降到0。


\textbf{出处：} 2021B*期中 Q13，原卷第4页（10分）


\textbf{答案依据：} 2021B*期中 Q13 随卷／配套参考答案；条件与解释经复核。
